\chapter{Linear Algebra}
\label{ch:linear_algebra}
% ============================================================================

Linear algebra is essential for state-space control, multi-input multi-output (MIMO) systems, and machine learning.

\section{Vectors and Matrices}

\subsubsection{Vector Operations}
\begin{itemize}
    \item Addition: $\vecx + \vecy = [x_1+y_1, x_2+y_2, \ldots, x_n+y_n]^\top$
    \item Scalar multiplication: $\alpha\vecx = [\alpha x_1, \alpha x_2, \ldots, \alpha x_n]^\top$
    \item Dot product: $\vecx \cdot \vecy = \vecx^\top\vecy = \sum_{i=1}^n x_i y_i$
    \item Norm: $\|\vecx\| = \sqrt{\vecx^\top\vecx}$ (Euclidean norm)
\end{itemize}

% Vector operations figure
\begin{figure}[htbp]
\centering
\begin{tikzpicture}

% ========== Panel (a): Vector Addition ==========
\begin{scope}[xshift=0cm]
    \node[font=\bfseries] at (2.5, 4.5) {(a) Vector Addition};

    % Axes
    \draw[->, thick] (-0.5, 0) -- (5.5, 0) node[right] {$x$};
    \draw[->, thick] (0, -0.5) -- (0, 4.2) node[above] {$y$};

    % Grid
    \draw[help lines, UofSC50Black!40] (0, 0) grid (5, 4);

    % Vector x (Garnet)
    \draw[->, very thick, UofSCGarnet] (0, 0) -- (3, 1);
    \node[UofSCGarnet, font=\bfseries, below right] at (1.5, 0.5) {$\vecx$};

    % Vector y (Atlantic) - starting from origin (for reference)
    \draw[->, very thick, UofSCAtlantic, dashed] (0, 0) -- (1, 2.5);
    \node[UofSCAtlantic, font=\bfseries, left] at (0.4, 1.3) {$\vecy$};

    % Vector y translated (Atlantic) - starting from tip of x
    \draw[->, very thick, UofSCAtlantic] (3, 1) -- (4, 3.5);
    \node[UofSCAtlantic, font=\bfseries, right] at (3.6, 2.3) {$\vecy$};

    % Resultant vector x + y (Horseshoe)
    \draw[->, line width=2pt, UofSCHorseshoe] (0, 0) -- (4, 3.5);
    \node[UofSCHorseshoe, font=\bfseries, above left] at (2.2, 2) {$\vecx + \vecy$};

    % Origin
    \node[below left, font=\small] at (0, 0) {$O$};

    % Parallelogram completion (dashed)
    \draw[dashed, UofSC70Black, thick] (1, 2.5) -- (4, 3.5);
\end{scope}

% ========== Panel (b): Scalar Multiplication ==========
\begin{scope}[xshift=8cm]
    \node[font=\bfseries] at (2.5, 4.5) {(b) Scalar Multiplication};

    % Axes
    \draw[->, thick] (-0.5, 0) -- (5.5, 0) node[right] {$x$};
    \draw[->, thick] (0, -0.5) -- (0, 4.2) node[above] {$y$};

    % Grid
    \draw[help lines, UofSC50Black!40] (0, 0) grid (5, 4);

    % Original vector v (Garnet)
    \draw[->, very thick, UofSCGarnet] (0, 0) -- (1.5, 1);
    \node[UofSCGarnet, font=\bfseries, below] at (1.2, 0.6) {$\vecv$};

    % 2v (Atlantic)
    \draw[->, very thick, UofSCAtlantic] (0, 0) -- (3, 2);
    \node[UofSCAtlantic, font=\bfseries, above left] at (2.8, 1.8) {$2\vecv$};

    % -v (Rose) - opposite direction
    \draw[->, very thick, UofSCRose] (0, 0) -- (-1.5, -1) node[below left, font=\bfseries] {$-\vecv$};

    % 0.5v (Horseshoe)
    \draw[->, very thick, UofSCHorseshoe] (0, 0) -- (0.75, 0.5);
    \node[UofSCHorseshoe, font=\bfseries, above] at (0.75, 0.5) {$\frac{1}{2}\vecv$};

    % Origin
    \node[below left, font=\small] at (0, 0) {$O$};

    % Annotation
    \node[align=left, font=\footnotesize, text width=3.5cm] at (4.2, 3.5) {
        Scalar $\alpha > 1$: stretch\\
        Scalar $0 < \alpha < 1$: shrink\\
        Scalar $\alpha < 0$: reverse
    };
\end{scope}

% ========== Panel (c): Dot Product ==========
\begin{scope}[xshift=4cm, yshift=-6cm]
    \node[font=\bfseries] at (2.5, 4.5) {(c) Dot Product: $\vecx \cdot \vecy = \|\vecx\| \|\vecy\| \cos\theta$};

    % Axes
    \draw[->, thick] (-0.5, 0) -- (5.5, 0) node[right] {$x$};
    \draw[->, thick] (0, -0.5) -- (0, 4.2) node[above] {$y$};

    % Grid
    \draw[help lines, UofSC50Black!40] (0, 0) grid (5, 4);

    % Vector x (Garnet)
    \draw[->, very thick, UofSCGarnet] (0, 0) -- (4.5, 0);
    \node[UofSCGarnet, font=\bfseries, below] at (2.25, -0.15) {$\vecx$};

    % Vector y (Atlantic)
    \draw[->, very thick, UofSCAtlantic] (0, 0) -- (3, 3);
    \node[UofSCAtlantic, font=\bfseries, above left] at (3, 3) {$\vecy$};

    % Angle arc (Horseshoe)
    \draw[->, thick, UofSCHorseshoe] (1, 0) arc (0:45:1);
    \node[UofSCHorseshoe, font=\small] at (1.4, 0.35) {$\theta$};

    % Projection of y onto x (Congaree)
    \draw[dashed, thick, UofSCCongaree] (3, 3) -- (3, 0);
    \node[circle, fill=UofSCCongaree, minimum size=5pt, inner sep=0pt] at (3, 0) {};

    % Projection segment
    \draw[line width=3pt, UofSCCongaree] (0, 0) -- (3, 0);
    \node[UofSCCongaree, font=\bfseries, below] at (1.5, -0.35) {$\|\vecy\|\cos\theta$};

    % Right angle marker
    \draw[thick, UofSCCongaree] (2.8, 0) -- (2.8, 0.2) -- (3, 0.2);

    % Origin
    \node[below left, font=\small] at (0, 0) {$O$};

    % Formula annotation
    \node[align=left, font=\footnotesize, text width=4cm, fill=white, draw=UofSC70Black] at (4.8, 2.8) {
        $\vecx \cdot \vecy = x_1 y_1 + x_2 y_2$\\[3pt]
        Geometric: projection of\\
        $\vecy$ onto $\vecx$, scaled by $\|\vecx\|$
    };
\end{scope}

\end{tikzpicture}
\caption{Geometric interpretation of fundamental vector operations in $\mathbb{R}^2$. (a) Vector addition follows the parallelogram law: $\vecx + \vecy$ is the diagonal of the parallelogram formed by $\vecx$ and $\vecy$. (b) Scalar multiplication scales vector magnitude; positive scalars preserve direction while negative scalars reverse it. (c) The dot product $\vecx \cdot \vecy = \|\vecx\|\|\vecy\|\cos\theta$ equals the product of $\|\vecx\|$ with the projection of $\vecy$ onto $\vecx$. The projection (Congaree) shows the geometric meaning of the inner product.}
\label{fig:vector_operations}
\end{figure}


\subsubsection{Matrix Operations}
\begin{itemize}
    \item Addition: $(\mA + \mB)_{ij} = a_{ij} + b_{ij}$
    \item Multiplication: $(\mA\mB)_{ij} = \sum_k a_{ik}b_{kj}$
    \item Transpose: $(\mA^\top)_{ij} = a_{ji}$
    \item Inverse: $\mA\mA^{-1} = \mA^{-1}\mA = \mI$ (if exists)
\end{itemize}

% Matrix transformation figure
\begin{figure}[htbp]
\centering
\begin{tikzpicture}

% ========== Panel (a): Original Unit Square ==========
\begin{scope}[xshift=0cm]
    \node[font=\bfseries] at (1.5, 3.8) {(a) Unit Square};

    % Axes
    \draw[->, thick] (-0.5, 0) -- (3.5, 0) node[right] {$x$};
    \draw[->, thick] (0, -0.5) -- (0, 3.5) node[above] {$y$};

    % Grid
    \draw[help lines, UofSC50Black!40] (-0.5, -0.5) grid (3, 3);

    % Unit square (filled)
    \fill[UofSCAtlantic!20] (0, 0) -- (1, 0) -- (1, 1) -- (0, 1) -- cycle;
    \draw[very thick, UofSCAtlantic] (0, 0) -- (1, 0) -- (1, 1) -- (0, 1) -- cycle;

    % Basis vectors
    \draw[->, line width=2pt, UofSCGarnet] (0, 0) -- (1, 0) node[below, font=\bfseries] {$\vece_1$};
    \draw[->, line width=2pt, UofSCHorseshoe] (0, 0) -- (0, 1) node[left, font=\bfseries] {$\vece_2$};

    % Corner labels
    \node[below left, font=\small] at (0, 0) {$O$};
    \node[circle, fill=UofSCBlack, minimum size=4pt, inner sep=0pt] at (0, 0) {};
    \node[circle, fill=UofSCBlack, minimum size=4pt, inner sep=0pt] at (1, 0) {};
    \node[circle, fill=UofSCBlack, minimum size=4pt, inner sep=0pt] at (1, 1) {};
    \node[circle, fill=UofSCBlack, minimum size=4pt, inner sep=0pt] at (0, 1) {};

    % Matrix label
    \node[align=center, font=\footnotesize] at (1.5, -1.2) {$\mI = \begin{bmatrix} 1 & 0 \\ 0 & 1 \end{bmatrix}$};
\end{scope}

% ========== Panel (b): Shear Transformation ==========
\begin{scope}[xshift=5cm]
    \node[font=\bfseries] at (2, 3.8) {(b) Shear};

    % Axes
    \draw[->, thick] (-0.5, 0) -- (4, 0) node[right] {$x$};
    \draw[->, thick] (0, -0.5) -- (0, 3.5) node[above] {$y$};

    % Grid
    \draw[help lines, UofSC50Black!40] (-0.5, -0.5) grid (3.5, 3);

    % Original unit square (ghost)
    \draw[thick, UofSC50Black, dashed] (0, 0) -- (1, 0) -- (1, 1) -- (0, 1) -- cycle;

    % Sheared parallelogram: A = [1 0.8; 0 1]
    % (0,0) -> (0,0), (1,0) -> (1,0), (0,1) -> (0.8,1), (1,1) -> (1.8,1)
    \fill[UofSCRose!20] (0, 0) -- (1, 0) -- (1.8, 1) -- (0.8, 1) -- cycle;
    \draw[very thick, UofSCRose] (0, 0) -- (1, 0) -- (1.8, 1) -- (0.8, 1) -- cycle;

    % Transformed basis vectors
    \draw[->, line width=2pt, UofSCGarnet] (0, 0) -- (1, 0) node[below, font=\bfseries] {$\mA\vece_1$};
    \draw[->, line width=2pt, UofSCHorseshoe] (0, 0) -- (0.8, 1) node[above left, font=\bfseries] {$\mA\vece_2$};

    % Corners
    \node[circle, fill=UofSCBlack, minimum size=4pt, inner sep=0pt] at (0, 0) {};
    \node[circle, fill=UofSCBlack, minimum size=4pt, inner sep=0pt] at (1, 0) {};
    \node[circle, fill=UofSCBlack, minimum size=4pt, inner sep=0pt] at (1.8, 1) {};
    \node[circle, fill=UofSCBlack, minimum size=4pt, inner sep=0pt] at (0.8, 1) {};

    % Matrix label
    \node[align=center, font=\footnotesize] at (2, -1.2) {$\mA = \begin{bmatrix} 1 & 0.8 \\ 0 & 1 \end{bmatrix}$};
\end{scope}

% ========== Panel (c): Rotation ==========
\begin{scope}[xshift=10cm]
    \node[font=\bfseries] at (1.5, 3.8) {(c) Rotation ($45^{\circ}$)};

    % Axes
    \draw[->, thick] (-1.5, 0) -- (2.5, 0) node[right] {$x$};
    \draw[->, thick] (0, -0.5) -- (0, 3.5) node[above] {$y$};

    % Grid
    \draw[help lines, UofSC50Black!40] (-1.5, -0.5) grid (2, 3);

    % Original unit square (ghost)
    \draw[thick, UofSC50Black, dashed] (0, 0) -- (1, 0) -- (1, 1) -- (0, 1) -- cycle;

    % Rotated square: R(45°) = [cos45 -sin45; sin45 cos45] = [0.707 -0.707; 0.707 0.707]
    % (0,0) -> (0,0), (1,0) -> (0.707, 0.707), (0,1) -> (-0.707, 0.707), (1,1) -> (0, 1.414)
    \fill[UofSCCongaree!20] (0, 0) -- (0.707, 0.707) -- (0, 1.414) -- (-0.707, 0.707) -- cycle;
    \draw[very thick, UofSCCongaree] (0, 0) -- (0.707, 0.707) -- (0, 1.414) -- (-0.707, 0.707) -- cycle;

    % Transformed basis vectors
    \draw[->, line width=2pt, UofSCGarnet] (0, 0) -- (0.707, 0.707) node[right, font=\bfseries] {$\mR\vece_1$};
    \draw[->, line width=2pt, UofSCHorseshoe] (0, 0) -- (-0.707, 0.707) node[left, font=\bfseries] {$\mR\vece_2$};

    % Rotation arc
    \draw[->, thick, UofSCAtlantic] (0.5, 0) arc (0:45:0.5);
    \node[UofSCAtlantic, font=\small] at (0.7, 0.25) {$45^{\circ}$};

    % Corners
    \node[circle, fill=UofSCBlack, minimum size=4pt, inner sep=0pt] at (0, 0) {};
    \node[circle, fill=UofSCBlack, minimum size=4pt, inner sep=0pt] at (0.707, 0.707) {};
    \node[circle, fill=UofSCBlack, minimum size=4pt, inner sep=0pt] at (0, 1.414) {};
    \node[circle, fill=UofSCBlack, minimum size=4pt, inner sep=0pt] at (-0.707, 0.707) {};

    % Matrix label
    \node[align=center, font=\footnotesize] at (0.5, -1.2) {$\mR = \begin{bmatrix} \cos 45^{\circ} & -\sin 45^{\circ} \\ \sin 45^{\circ} & \cos 45^{\circ} \end{bmatrix}$};
\end{scope}

% ========== Panel (d): Scaling ==========
\begin{scope}[xshift=0cm, yshift=-6.5cm]
    \node[font=\bfseries] at (2.5, 3.8) {(d) Non-uniform Scaling};

    % Axes
    \draw[->, thick] (-0.5, 0) -- (5, 0) node[right] {$x$};
    \draw[->, thick] (0, -0.5) -- (0, 3.5) node[above] {$y$};

    % Grid
    \draw[help lines, UofSC50Black!40] (-0.5, -0.5) grid (4.5, 3);

    % Original unit square (ghost)
    \draw[thick, UofSC50Black, dashed] (0, 0) -- (1, 0) -- (1, 1) -- (0, 1) -- cycle;

    % Scaled rectangle: S = [2 0; 0 0.5]
    % (0,0) -> (0,0), (1,0) -> (2,0), (0,1) -> (0,0.5), (1,1) -> (2,0.5)
    \fill[UofSCHorseshoe!20] (0, 0) -- (2, 0) -- (2, 0.5) -- (0, 0.5) -- cycle;
    \draw[very thick, UofSCHorseshoe] (0, 0) -- (2, 0) -- (2, 0.5) -- (0, 0.5) -- cycle;

    % Transformed basis vectors
    \draw[->, line width=2pt, UofSCGarnet] (0, 0) -- (2, 0) node[below, font=\bfseries] {$\mS\vece_1$};
    \draw[->, line width=2pt, UofSCAtlantic] (0, 0) -- (0, 0.5) node[left, font=\bfseries] {$\mS\vece_2$};

    % Corners
    \node[circle, fill=UofSCBlack, minimum size=4pt, inner sep=0pt] at (0, 0) {};
    \node[circle, fill=UofSCBlack, minimum size=4pt, inner sep=0pt] at (2, 0) {};
    \node[circle, fill=UofSCBlack, minimum size=4pt, inner sep=0pt] at (2, 0.5) {};
    \node[circle, fill=UofSCBlack, minimum size=4pt, inner sep=0pt] at (0, 0.5) {};

    % Annotations
    \node[UofSCGarnet, font=\footnotesize] at (3.2, 0.15) {stretch by 2};
    \node[UofSCAtlantic, font=\footnotesize] at (1.8, 1.2) {shrink by $\frac{1}{2}$};

    % Matrix label
    \node[align=center, font=\footnotesize] at (2.5, -1.2) {$\mS = \begin{bmatrix} 2 & 0 \\ 0 & 0.5 \end{bmatrix}$};
\end{scope}

% ========== Panel (e): Combined Transformation ==========
\begin{scope}[xshift=7cm, yshift=-6.5cm]
    \node[font=\bfseries] at (3, 3.8) {(e) Combined: Rotation + Scaling};

    % Axes
    \draw[->, thick] (-1.5, 0) -- (5, 0) node[right] {$x$};
    \draw[->, thick] (0, -1.5) -- (0, 3.5) node[above] {$y$};

    % Grid
    \draw[help lines, UofSC50Black!40] (-1.5, -1.5) grid (4.5, 3);

    % Original unit square (ghost)
    \draw[thick, UofSC50Black, dashed] (0, 0) -- (1, 0) -- (1, 1) -- (0, 1) -- cycle;

    % Combined: Scale by [1.5, 0.8] then rotate 30°
    % First scale: (1,0) -> (1.5, 0), (0,1) -> (0, 0.8)
    % Then rotate 30°: [cos30 -sin30; sin30 cos30]
    % R*S*e1 = R*(1.5, 0) = (1.5*cos30, 1.5*sin30) = (1.299, 0.75)
    % R*S*e2 = R*(0, 0.8) = (-0.8*sin30, 0.8*cos30) = (-0.4, 0.693)
    % (1,1) -> R*S*(1,1) = R*(1.5, 0.8) = (1.5*cos30 - 0.8*sin30, 1.5*sin30 + 0.8*cos30) = (0.899, 1.443)
    \fill[UofSCGarnet!20] (0, 0) -- (1.299, 0.75) -- (0.899, 1.443) -- (-0.4, 0.693) -- cycle;
    \draw[very thick, UofSCGarnet] (0, 0) -- (1.299, 0.75) -- (0.899, 1.443) -- (-0.4, 0.693) -- cycle;

    % Transformed basis vectors
    \draw[->, line width=2pt, UofSCAtlantic] (0, 0) -- (1.299, 0.75) node[right, font=\bfseries] {$\mM\vece_1$};
    \draw[->, line width=2pt, UofSCHorseshoe] (0, 0) -- (-0.4, 0.693) node[above left, font=\bfseries] {$\mM\vece_2$};

    % Corners
    \node[circle, fill=UofSCBlack, minimum size=4pt, inner sep=0pt] at (0, 0) {};
    \node[circle, fill=UofSCBlack, minimum size=4pt, inner sep=0pt] at (1.299, 0.75) {};
    \node[circle, fill=UofSCBlack, minimum size=4pt, inner sep=0pt] at (0.899, 1.443) {};
    \node[circle, fill=UofSCBlack, minimum size=4pt, inner sep=0pt] at (-0.4, 0.693) {};

    % Matrix label
    \node[align=center, font=\footnotesize] at (3, -1.2) {$\mM = \mR(30^{\circ}) \cdot \mS = \begin{bmatrix} 1.30 & -0.40 \\ 0.75 & 0.69 \end{bmatrix}$};
\end{scope}

\end{tikzpicture}
\caption{Geometric interpretation of $2 \times 2$ matrix transformations acting on the unit square. (a) The identity matrix preserves all vectors. (b) Shear transforms the square into a parallelogram by sliding points parallel to one axis. (c) Rotation preserves lengths and angles (orthogonal transformation). (d) Non-uniform scaling stretches or compresses along each axis independently. (e) Combined transformations compose: here scaling followed by rotation. In each case, the matrix columns represent the transformed basis vectors $\mA\vece_1$ and $\mA\vece_2$. The dashed squares show the original unit square for reference.}
\label{fig:matrix_transformation}
\end{figure}


\section{Special Matrices}

\begin{itemize}
    \item \textbf{Identity matrix:} $\mI$ with ones on diagonal, zeros elsewhere
    \item \textbf{Diagonal matrix:} $\mD = \diag(d_1, d_2, \ldots, d_n)$
    \item \textbf{Symmetric matrix:} $\mA = \mA^\top$
    \item \textbf{Positive definite:} $\vecx^\top\mA\vecx > 0$ for all $\vecx \neq \mathbf{0}$
    \item \textbf{Orthogonal matrix:} $\mQ^\top\mQ = \mQ\mQ^\top = \mI$
\end{itemize}

\section{Determinant and Rank}

The \textbf{determinant} of a square matrix:
\begin{equation}
\det(\mA) = |\mA| = \sum_{j=1}^n (-1)^{1+j} a_{1j} M_{1j}
\end{equation}
where $M_{1j}$ is the minor (determinant of submatrix).

Properties:
\begin{itemize}
    \item $\det(\mA\mB) = \det(\mA)\det(\mB)$
    \item $\det(\mA^{-1}) = 1/\det(\mA)$
    \item $\det(\mA^\top) = \det(\mA)$
    \item $\mA$ is invertible if and only if $\det(\mA) \neq 0$
\end{itemize}

The \textbf{rank} is the number of linearly independent rows (or columns):
\begin{equation}
\rank(\mA) \leq \min(m, n) \text{ for } \mA \in \mathbb{R}^{m \times n}
\end{equation}

\section{Eigenvalues and Eigenvectors}

\begin{definitionbox}[title=Definition 0.5: Eigenvalue Problem]
For a square matrix $\mA$, scalar $\lambda$ is an \textbf{eigenvalue} and non-zero vector $\vecv$ is the corresponding \textbf{eigenvector} if:
\begin{equation}
\mA\vecv = \lambda\vecv
\end{equation}
\end{definitionbox}

Eigenvalues are roots of the \textbf{characteristic equation}:
\begin{equation}
\det(\mA - \lambda\mI) = 0
\end{equation}

% Eigenvalue visualization figure
\begin{figure}[htbp]
\centering
\begin{tikzpicture}

% ========== Main Panel: Eigenvector Visualization ==========
\begin{scope}[xshift=0cm]
    \node[font=\bfseries\large] at (4, 5.5) {Eigenvectors: Directions Preserved Under Transformation};

    % Axes
    \draw[->, thick] (-3.5, 0) -- (5.5, 0) node[right] {$x$};
    \draw[->, thick] (0, -3.5) -- (0, 5) node[above] {$y$};

    % Grid
    \draw[help lines, UofSC50Black!30] (-3, -3) grid (5, 4.5);

    % Matrix A = [2 1; 1 2] has eigenvalues 3 and 1
    % Eigenvector for lambda=3: v1 = [1; 1] (direction 45 degrees)
    % Eigenvector for lambda=1: v2 = [1; -1] (direction -45 degrees)

    % ========== Eigenvector 1: lambda = 3 ==========
    % Original eigenvector v1 (shorter, dashed)
    \draw[->, line width=2pt, UofSCGarnet, dashed] (0, 0) -- (1.2, 1.2);
    \node[UofSCGarnet, font=\bfseries, above left] at (0.9, 0.9) {$\vecv_1$};

    % Transformed Av1 = 3*v1 (longer, solid)
    \draw[->, line width=2.5pt, UofSCGarnet] (0, 0) -- (3.6, 3.6);
    \node[UofSCGarnet, font=\bfseries, above right] at (3.4, 3.4) {$\mA\vecv_1 = 3\vecv_1$};

    % Scaling annotation for v1
    \draw[<->, UofSC70Black, thick] (1.4, 1.0) -- (3.4, 3.0);
    \node[UofSC70Black, font=\footnotesize, fill=white, inner sep=1pt] at (2.4, 2.0) {$\times 3$};

    % ========== Eigenvector 2: lambda = 1 ==========
    % Original eigenvector v2 (dashed)
    \draw[->, line width=2pt, UofSCAtlantic, dashed] (0, 0) -- (1.5, -1.5);
    \node[UofSCAtlantic, font=\bfseries, below right] at (1.2, -1.2) {$\vecv_2$};

    % Transformed Av2 = 1*v2 (same length, solid - overlaps)
    \draw[->, line width=2.5pt, UofSCAtlantic] (0, 0) -- (1.5, -1.5);
    \node[UofSCAtlantic, font=\bfseries, below left] at (0.6, -1.8) {$\mA\vecv_2 = 1 \cdot \vecv_2$};

    % Scaling annotation for v2
    \node[UofSC70Black, font=\footnotesize, fill=white, inner sep=2pt] at (2.5, -1.5) {$\times 1$ (unchanged)};

    % ========== Non-eigenvector for comparison ==========
    % Original vector w (not an eigenvector)
    \draw[->, line width=1.5pt, UofSCHorseshoe, dashed] (0, 0) -- (2, 0.5);
    \node[UofSCHorseshoe, font=\bfseries, above] at (1.8, 0.5) {$\vecw$};

    % Transformed Aw (direction changes!)
    % A*[2; 0.5] = [2*2 + 1*0.5; 1*2 + 2*0.5] = [4.5; 3]
    \draw[->, line width=2pt, UofSCHorseshoe] (0, 0) -- (4.5, 3);
    \node[UofSCHorseshoe, font=\bfseries, above left] at (4.3, 2.8) {$\mA\vecw$};

    % Direction change annotation
    \draw[->, UofSCRose, thick, dashed] (2.2, 0.7) to[bend left=20] (3.8, 2.2);
    \node[UofSCRose, font=\footnotesize, fill=white, inner sep=1pt] at (3.5, 1.2) {direction changes!};

    % Origin
    \node[below left, font=\small] at (0, 0) {$O$};
    \node[circle, fill=UofSCBlack, minimum size=5pt, inner sep=0pt] at (0, 0) {};

    % ========== Matrix and Eigenvalue Info Box ==========
    \node[draw=UofSC70Black, thick, fill=white, align=left, font=\footnotesize,
          text width=5.5cm, inner sep=8pt] at (8, 2) {
        \textbf{Matrix:}\\[2pt]
        $\mA = \begin{bmatrix} 2 & 1 \\ 1 & 2 \end{bmatrix}$\\[8pt]
        \textbf{Characteristic equation:}\\[2pt]
        $\det(\mA - \lambda\mI) = \lambda^2 - 4\lambda + 3 = 0$\\[8pt]
        \textbf{Eigenvalues:}\\[2pt]
        $\lambda_1 = 3$, \quad $\lambda_2 = 1$\\[8pt]
        \textbf{Eigenvectors:}\\[2pt]
        $\vecv_1 = \begin{bmatrix} 1 \\ 1 \end{bmatrix}$ for $\lambda_1 = 3$\\[4pt]
        $\vecv_2 = \begin{bmatrix} 1 \\ -1 \end{bmatrix}$ for $\lambda_2 = 1$
    };

    % ========== Key Concept Box ==========
    \node[draw=UofSCGarnet, thick, fill=UofSCGarnet!8, align=left, font=\footnotesize,
          text width=5.5cm, inner sep=8pt] at (8, -2) {
        \textbf{Key Insight:}\\[4pt]
        Eigenvectors are special directions that remain\\
        \textit{parallel to themselves} under transformation.\\[4pt]
        The eigenvalue $\lambda$ is the \textbf{scaling factor}:\\[2pt]
        \textbullet\ $|\lambda| > 1$: stretching\\
        \textbullet\ $|\lambda| < 1$: compression\\
        \textbullet\ $\lambda < 0$: reversal\\[4pt]
        Non-eigenvectors (like $\vecw$) change direction.
    };

\end{scope}

\end{tikzpicture}
\caption{Geometric visualization of eigenvalues and eigenvectors. For matrix $\mA$, eigenvectors $\vecv_1$ and $\vecv_2$ (dashed) are special directions that remain parallel to themselves after transformation (solid). The eigenvalue is the scaling factor: $\vecv_1$ is scaled by $\lambda_1 = 3$ (stretched), while $\vecv_2$ is scaled by $\lambda_2 = 1$ (unchanged). In contrast, a non-eigenvector $\vecw$ both rotates and scales under transformation---its direction is not preserved. This geometric property is fundamental to understanding system stability: eigenvalues determine how state components grow or decay along eigenvector directions.}
\label{fig:eigenvalue_visualization}
\end{figure}


\subsubsection{Properties}
\begin{itemize}
    \item Sum of eigenvalues = $\trace(\mA) = \sum_i a_{ii}$
    \item Product of eigenvalues = $\det(\mA)$
    \item Eigenvalues of $\mA^{-1}$ are $1/\lambda_i$
    \item Eigenvalues of $\mA^n$ are $\lambda_i^n$
\end{itemize}

\begin{keyconceptbox}[title=Stability and Eigenvalues]
A continuous-time system $\dot{\vecx} = \mA\vecx$ is:
\begin{itemize}
    \item \textbf{Stable} if all eigenvalues of $\mA$ have negative real parts
    \item \textbf{Marginally stable} if eigenvalues have non-positive real parts with simple imaginary-axis eigenvalues
    \item \textbf{Unstable} if any eigenvalue has positive real part
\end{itemize}
\end{keyconceptbox}

\begin{examplebox}[title=Example: Eigenvalues and Eigenvectors of a 2x2 Matrix]
Find the eigenvalues and eigenvectors of:
\[
\mA = \begin{bmatrix} 4 & 1 \\ 2 & 3 \end{bmatrix}
\]

\textbf{Solution:}

\textbf{Step 1: Find eigenvalues} by solving $\det(\mA - \lambda\mI) = 0$:
\[
\det\begin{bmatrix} 4-\lambda & 1 \\ 2 & 3-\lambda \end{bmatrix} = (4-\lambda)(3-\lambda) - 2 = 0
\]
\[
\lambda^2 - 7\lambda + 12 - 2 = \lambda^2 - 7\lambda + 10 = 0
\]
\[
(\lambda - 5)(\lambda - 2) = 0 \quad \Rightarrow \quad \lambda_1 = 5, \quad \lambda_2 = 2
\]

\textbf{Step 2: Find eigenvector for $\lambda_1 = 5$}:
\[
(\mA - 5\mI)\vecv_1 = \begin{bmatrix} -1 & 1 \\ 2 & -2 \end{bmatrix}\vecv_1 = \mathbf{0}
\]
From $-v_1 + v_2 = 0$, we get $v_2 = v_1$. Choosing $v_1 = 1$:
\[
\vecv_1 = \begin{bmatrix} 1 \\ 1 \end{bmatrix}
\]

\textbf{Step 3: Find eigenvector for $\lambda_2 = 2$}:
\[
(\mA - 2\mI)\vecv_2 = \begin{bmatrix} 2 & 1 \\ 2 & 1 \end{bmatrix}\vecv_2 = \mathbf{0}
\]
From $2v_1 + v_2 = 0$, we get $v_2 = -2v_1$. Choosing $v_1 = 1$:
\[
\vecv_2 = \begin{bmatrix} 1 \\ -2 \end{bmatrix}
\]

\textbf{Verification:} $\mA\vecv_1 = \begin{bmatrix} 5 \\ 5 \end{bmatrix} = 5\vecv_1$ \checkmark \quad and \quad $\mA\vecv_2 = \begin{bmatrix} 2 \\ -4 \end{bmatrix} = 2\vecv_2$ \checkmark
\end{examplebox}

\begin{keyconceptbox}[title=Cayley-Hamilton Theorem]
Every square matrix satisfies its own characteristic equation. If the characteristic polynomial of $\mA$ is:
\[
p(\lambda) = \det(\lambda\mI - \mA) = \lambda^n + a_{n-1}\lambda^{n-1} + \cdots + a_1\lambda + a_0
\]
then:
\[
p(\mA) = \mA^n + a_{n-1}\mA^{n-1} + \cdots + a_1\mA + a_0\mI = \mathbf{0}
\]

\textbf{Applications:}
\begin{itemize}
    \item Computing $\mA^{-1}$ from powers of $\mA$
    \item Expressing high powers of $\mA$ in terms of lower powers
    \item Computing the matrix exponential $\ee^{\mA t}$
\end{itemize}
\end{keyconceptbox}

\section{Matrix Decompositions}

\subsubsection{Eigendecomposition}
For diagonalizable matrix $\mA$:
\begin{equation}
\mA = \mV\boldsymbol{\Lambda}\mV^{-1}
\end{equation}
where $\mV$ contains eigenvectors and $\boldsymbol{\Lambda} = \diag(\lambda_1, \ldots, \lambda_n)$.

\subsubsection{Singular Value Decomposition (SVD)}
Any matrix $\mA \in \mathbb{R}^{m \times n}$:
\begin{equation}
\mA = \mU\boldsymbol{\Sigma}\mV^\top
\end{equation}
where $\mU, \mV$ are orthogonal and $\boldsymbol{\Sigma}$ contains singular values.

\begin{examplebox}[title=Example: Computing SVD for a Simple Matrix]
Find the SVD of:
\[
\mA = \begin{bmatrix} 3 & 0 \\ 0 & 2 \end{bmatrix}
\]

\textbf{Solution:}

\textbf{Step 1: Compute $\mA^\top\mA$}:
\[
\mA^\top\mA = \begin{bmatrix} 3 & 0 \\ 0 & 2 \end{bmatrix}\begin{bmatrix} 3 & 0 \\ 0 & 2 \end{bmatrix} = \begin{bmatrix} 9 & 0 \\ 0 & 4 \end{bmatrix}
\]

\textbf{Step 2: Find eigenvalues of $\mA^\top\mA$} (these are $\sigma_i^2$):

Since this is diagonal, eigenvalues are $\sigma_1^2 = 9$ and $\sigma_2^2 = 4$.

Singular values: $\sigma_1 = 3$, $\sigma_2 = 2$.

\textbf{Step 3: Form $\boldsymbol{\Sigma}$}:
\[
\boldsymbol{\Sigma} = \begin{bmatrix} 3 & 0 \\ 0 & 2 \end{bmatrix}
\]

\textbf{Step 4: Find $\mV$} (eigenvectors of $\mA^\top\mA$):

For diagonal matrix, eigenvectors are standard basis vectors:
\[
\mV = \begin{bmatrix} 1 & 0 \\ 0 & 1 \end{bmatrix} = \mI
\]

\textbf{Step 5: Find $\mU$} from $\mU = \mA\mV\boldsymbol{\Sigma}^{-1}$:
\[
\mU = \begin{bmatrix} 3 & 0 \\ 0 & 2 \end{bmatrix}\begin{bmatrix} 1 & 0 \\ 0 & 1 \end{bmatrix}\begin{bmatrix} 1/3 & 0 \\ 0 & 1/2 \end{bmatrix} = \begin{bmatrix} 1 & 0 \\ 0 & 1 \end{bmatrix} = \mI
\]

\textbf{Result:} $\mA = \mU\boldsymbol{\Sigma}\mV^\top = \mI \cdot \boldsymbol{\Sigma} \cdot \mI = \boldsymbol{\Sigma}$ (as expected for a diagonal matrix).

For non-diagonal matrices, the procedure involves more computation but follows the same principles.
\end{examplebox}

\subsubsection{QR Decomposition}
\begin{equation}
\mA = \mQ\mR
\end{equation}
where $\mQ$ is orthogonal and $\mR$ is upper triangular.

\subsubsection{LU Decomposition}
\begin{equation}
\mA = \mL\mU
\end{equation}
where $\mL$ is lower triangular and $\mU$ is upper triangular.

\section{Matrix Exponential}

The matrix exponential is crucial for solving state-space equations:

\begin{definitionbox}[title=Definition 0.6: Matrix Exponential]
\begin{equation}
\ee^{\mA t} = \mI + \mA t + \frac{(\mA t)^2}{2!} + \frac{(\mA t)^3}{3!} + \cdots = \sum_{k=0}^{\infty} \frac{(\mA t)^k}{k!}
\end{equation}
\end{definitionbox}

% Matrix exponential figure
\begin{figure}[htbp]
\centering
\begin{tikzpicture}

% ========== Panel (a): Stable Spiral (Complex eigenvalues with negative real part) ==========
\begin{scope}[xshift=0cm, yshift=0cm]
    \node[font=\bfseries] at (3.5, 4.5) {(a) Stable Spiral};
    \node[font=\footnotesize, UofSCHorseshoe] at (3.5, 4.0) {$\lambda = -0.3 \pm 0.8\jj$};

    % Use TikZ directly for the spiral instead of pgfplots
    \begin{scope}[xshift=0.5cm, yshift=0.5cm, scale=1.3]
        % Grid and axes
        \draw[help lines, gray!20] (-2,-2) grid (2,2);
        \draw[thick, ->, >=stealth] (-2.3,0) -- (2.3,0) node[right, font=\small] {$x_1$};
        \draw[thick, ->, >=stealth] (0,-2.3) -- (0,2.3) node[above, font=\small] {$x_2$};

        % Spiral trajectory 1 - Garnet (using parametric with radians)
        \draw[UofSCGarnet, very thick, smooth, samples=100, domain=0:6.28]
            plot ({2*exp(-0.25*\x)*cos(\x r)}, {2*exp(-0.25*\x)*sin(\x r)});
        \draw[UofSCGarnet, thick, ->] (1.2, 0.8) -- (1.1, 0.9);

        % Spiral trajectory 2 - Atlantic
        \draw[UofSCAtlantic, very thick, smooth, samples=100, domain=0:6.28]
            plot ({1.8*exp(-0.25*\x)*cos(\x r + 2.09)}, {1.8*exp(-0.25*\x)*sin(\x r + 2.09)});
        \draw[UofSCAtlantic, thick, ->] (-0.6, 1.2) -- (-0.7, 1.15);

        % Spiral trajectory 3 - Horseshoe
        \draw[UofSCHorseshoe, very thick, smooth, samples=100, domain=0:6.28]
            plot ({1.6*exp(-0.25*\x)*cos(\x r + 4.19)}, {1.6*exp(-0.25*\x)*sin(\x r + 4.19)});
        \draw[UofSCHorseshoe, thick, ->] (-0.5, -0.9) -- (-0.4, -0.95);

        % Equilibrium point
        \node[circle, fill=UofSCBlack, minimum size=5pt, inner sep=0pt] at (0,0) {};

        % Initial condition markers
        \node[circle, fill=UofSCGarnet, minimum size=4pt, inner sep=0pt] at (2,0) {};
        \node[circle, fill=UofSCAtlantic, minimum size=4pt, inner sep=0pt] at (-0.9,1.56) {};
        \node[circle, fill=UofSCHorseshoe, minimum size=4pt, inner sep=0pt] at (-0.8,-1.39) {};
    \end{scope}

    % s-plane inset
    \begin{scope}[xshift=5.5cm, yshift=0.5cm, scale=0.4]
        \fill[UofSCHorseshoe!15] (-2,-2) rectangle (0,2);
        \draw[->, thick] (-2.2,0) -- (1,0) node[right, font=\tiny] {$\Real$};
        \draw[->, thick] (0,-2.2) -- (0,2.2) node[above, font=\tiny] {$\Imag$};
        \node[circle, fill=UofSCGarnet, minimum size=5pt, inner sep=0pt] at (-0.6,1.6) {};
        \node[circle, fill=UofSCGarnet, minimum size=5pt, inner sep=0pt] at (-0.6,-1.6) {};
        \draw[thick, UofSCCongaree] (0,-2) -- (0,2);
        \node[font=\tiny, UofSCHorseshoe] at (-1.2,-1.8) {LHP};
    \end{scope}
\end{scope}

% ========== Panel (b): Unstable Node (Real positive eigenvalues) ==========
\begin{scope}[xshift=8cm, yshift=0cm]
    \node[font=\bfseries] at (3.5, 4.5) {(b) Unstable Node};
    \node[font=\footnotesize, UofSCRose] at (3.5, 4.0) {$\lambda_1 = 0.5$, $\lambda_2 = 1.2$};

    % Use TikZ directly
    \begin{scope}[xshift=0.5cm, yshift=0.5cm, scale=1.3]
        % Grid and axes
        \draw[help lines, gray!20] (-2,-2) grid (2,2);
        \draw[thick, ->, >=stealth] (-2.3,0) -- (2.3,0) node[right, font=\small] {$x_1$};
        \draw[thick, ->, >=stealth] (0,-2.3) -- (0,2.3) node[above, font=\small] {$x_2$};

        % Trajectory 1 - Garnet (outward from origin)
        \draw[UofSCGarnet, very thick, smooth, samples=50, domain=0:1.4]
            plot ({0.3*exp(0.5*\x)}, {0.2*exp(1.2*\x)});
        \draw[UofSCGarnet, thick, ->] (0.4, 0.5) -- (0.45, 0.7);

        % Trajectory 2 - Atlantic
        \draw[UofSCAtlantic, very thick, smooth, samples=50, domain=0:1.4]
            plot ({-0.25*exp(0.5*\x)}, {0.15*exp(1.2*\x)});
        \draw[UofSCAtlantic, thick, ->] (-0.35, 0.4) -- (-0.38, 0.55);

        % Trajectory 3 - Horseshoe
        \draw[UofSCHorseshoe, very thick, smooth, samples=50, domain=0:1.4]
            plot ({0.2*exp(0.5*\x)}, {-0.3*exp(1.2*\x)});
        \draw[UofSCHorseshoe, thick, ->] (0.3, -0.7) -- (0.32, -0.9);

        % Trajectory 4 - Rose
        \draw[UofSCRose, very thick, smooth, samples=50, domain=0:1.4]
            plot ({-0.35*exp(0.5*\x)}, {-0.2*exp(1.2*\x)});
        \draw[UofSCRose, thick, ->] (-0.5, -0.5) -- (-0.55, -0.7);

        % Equilibrium point
        \node[circle, fill=UofSCBlack, minimum size=5pt, inner sep=0pt] at (0,0) {};

        % Initial condition markers
        \node[circle, fill=UofSCGarnet, minimum size=4pt, inner sep=0pt] at (0.3,0.2) {};
        \node[circle, fill=UofSCAtlantic, minimum size=4pt, inner sep=0pt] at (-0.25,0.15) {};
        \node[circle, fill=UofSCHorseshoe, minimum size=4pt, inner sep=0pt] at (0.2,-0.3) {};
        \node[circle, fill=UofSCRose, minimum size=4pt, inner sep=0pt] at (-0.35,-0.2) {};
    \end{scope}

    % s-plane inset
    \begin{scope}[xshift=5.5cm, yshift=0.5cm, scale=0.4]
        \fill[UofSCRose!15] (0,-2) rectangle (2,2);
        \draw[->, thick] (-1,0) -- (2.2,0) node[right, font=\tiny] {$\Real$};
        \draw[->, thick] (0,-2.2) -- (0,2.2) node[above, font=\tiny] {$\Imag$};
        \node[circle, fill=UofSCGarnet, minimum size=5pt, inner sep=0pt] at (1,0) {};
        \node[circle, fill=UofSCGarnet, minimum size=5pt, inner sep=0pt] at (2,0) {};
        \draw[thick, UofSCCongaree] (0,-2) -- (0,2);
        \node[font=\tiny, UofSCRose] at (1.2,-1.8) {RHP};
    \end{scope}
\end{scope}

% ========== Panel (c): Saddle Point (One positive, one negative real eigenvalue) ==========
\begin{scope}[xshift=4cm, yshift=-8cm]
    \node[font=\bfseries] at (3.5, 4.5) {(c) Saddle Point};
    \node[font=\footnotesize, UofSCCongaree] at (3.5, 4.0) {$\lambda_1 = -1$, $\lambda_2 = 0.8$};

    % Use TikZ directly
    \begin{scope}[xshift=0.5cm, yshift=0.5cm, scale=1.3]
        % Grid and axes
        \draw[help lines, gray!20] (-2,-2) grid (2,2);
        \draw[thick, ->, >=stealth] (-2.3,0) -- (2.3,0) node[right, font=\small] {$x_1$};
        \draw[thick, ->, >=stealth] (0,-2.3) -- (0,2.3) node[above, font=\small] {$x_2$};

        % Stable manifold (along x1 axis, lambda = -1)
        \draw[UofSCHorseshoe, very thick, dashed] (-2.2,0) -- (2.2,0);
        \node[UofSCHorseshoe, font=\footnotesize] at (2,0.25) {stable};

        % Unstable manifold (along x2 axis, lambda = 0.8)
        \draw[UofSCRose, very thick, dashed] (0,-2.2) -- (0,2.2);
        \node[UofSCRose, font=\footnotesize, rotate=90] at (0.3,1.8) {unstable};

        % Saddle trajectories - hyperbolic curves
        \draw[UofSCGarnet, very thick, smooth, samples=50, domain=-1:1]
            plot ({1.5*exp(-\x)}, {0.3*exp(0.8*\x)});
        \draw[UofSCGarnet, thick, ->] (0.8, 0.4) -- (0.65, 0.5);

        \draw[UofSCAtlantic, very thick, smooth, samples=50, domain=-1:1]
            plot ({-1.5*exp(-\x)}, {0.25*exp(0.8*\x)});
        \draw[UofSCAtlantic, thick, ->] (-0.7, 0.35) -- (-0.55, 0.45);

        \draw[UofSCCongaree, very thick, smooth, samples=50, domain=-1:1]
            plot ({1.2*exp(-\x)}, {-0.35*exp(0.8*\x)});
        \draw[UofSCCongaree, thick, ->] (0.65, -0.45) -- (0.5, -0.55);

        \draw[UofSCRose, very thick, smooth, samples=50, domain=-1:1]
            plot ({-1.3*exp(-\x)}, {-0.28*exp(0.8*\x)});
        \draw[UofSCRose, thick, ->] (-0.7, -0.35) -- (-0.55, -0.45);

        % Equilibrium point
        \node[circle, fill=UofSCBlack, minimum size=5pt, inner sep=0pt] at (0,0) {};
    \end{scope}

    % s-plane inset
    \begin{scope}[xshift=5.5cm, yshift=0.5cm, scale=0.4]
        \fill[UofSCHorseshoe!15] (-2,-2) rectangle (0,2);
        \fill[UofSCRose!15] (0,-2) rectangle (2,2);
        \draw[->, thick] (-2.2,0) -- (2.2,0) node[right, font=\tiny] {$\Real$};
        \draw[->, thick] (0,-2.2) -- (0,2.2) node[above, font=\tiny] {$\Imag$};
        \node[circle, fill=UofSCGarnet, minimum size=5pt, inner sep=0pt] at (-2,0) {};
        \node[circle, fill=UofSCGarnet, minimum size=5pt, inner sep=0pt] at (1.6,0) {};
        \draw[thick, UofSCCongaree] (0,-2) -- (0,2);
    \end{scope}
\end{scope}

% ========== Legend/Key Box ==========
\node[draw=UofSC70Black, thick, fill=white, align=left, font=\footnotesize,
      text width=7cm, inner sep=10pt] at (12.5, -2) {
    \textbf{State Evolution: $\vecx(t) = \ee^{\mA t}\vecx_0$}\\[6pt]
    \textbf{Stability Classification:}\\[4pt]
    \textbullet\ \textbf{Stable spiral:} $\Real(\lambda) < 0$, complex pair\\
    \quad Decaying oscillations toward equilibrium\\[3pt]
    \textbullet\ \textbf{Unstable node:} $\lambda_1, \lambda_2 > 0$, real\\
    \quad All trajectories diverge from origin\\[3pt]
    \textbullet\ \textbf{Saddle point:} $\lambda_1 < 0 < \lambda_2$, real\\
    \quad Stable/unstable manifolds, overall unstable\\[6pt]
    \textbf{Insets:} Eigenvalue locations in $s$-plane\\
    LHP (green) = stable, RHP (red) = unstable
};

\end{tikzpicture}
\caption{Phase portraits showing time evolution of the autonomous system $\dot{\vecx} = \mA\vecx$ via the matrix exponential $\vecx(t) = \ee^{\mA t}\vecx_0$ for three characteristic eigenvalue configurations. (a) \textbf{Stable spiral}: complex conjugate eigenvalues with negative real part produce decaying oscillatory trajectories spiraling toward the origin. (b) \textbf{Unstable node}: two positive real eigenvalues cause all trajectories to diverge exponentially from the equilibrium. (c) \textbf{Saddle point}: one stable and one unstable eigenvalue create invariant manifolds; trajectories approach the stable manifold but ultimately diverge along the unstable direction. Insets show eigenvalue locations in the $s$-plane; the imaginary axis separates stable (LHP) from unstable (RHP) behavior.}
\label{fig:matrix_exponential}
\end{figure}


\subsubsection{Properties}
\begin{itemize}
    \item $\ee^{\mathbf{0}} = \mI$
    \item $\frac{\dd}{\dd t}\ee^{\mA t} = \mA\ee^{\mA t}$
    \item $\ee^{\mA(t_1+t_2)} = \ee^{\mA t_1}\ee^{\mA t_2}$
    \item $(\ee^{\mA t})^{-1} = \ee^{-\mA t}$
\end{itemize}

\subsubsection{Computation via Laplace}
\begin{equation}
\ee^{\mA t} = \iLaplace\{(s\mI - \mA)^{-1}\}
\end{equation}

\subsubsection{Computation via eigendecomposition}
If $\mA = \mV\boldsymbol{\Lambda}\mV^{-1}$:
\begin{equation}
\ee^{\mA t} = \mV\ee^{\boldsymbol{\Lambda} t}\mV^{-1} = \mV\diag(\ee^{\lambda_1 t}, \ldots, \ee^{\lambda_n t})\mV^{-1}
\end{equation}

\begin{examplebox}[title=Example: Computing Matrix Exponential for a Simple System]
Compute $\ee^{\mA t}$ for the system matrix:
\[
\mA = \begin{bmatrix} -1 & 0 \\ 0 & -2 \end{bmatrix}
\]

\textbf{Solution:}

\textbf{Method 1: Direct computation (diagonal matrix)}

For a diagonal matrix, the matrix exponential is simply:
\[
\ee^{\mA t} = \begin{bmatrix} \ee^{-t} & 0 \\ 0 & \ee^{-2t} \end{bmatrix}
\]

\textbf{Method 2: Using the Laplace transform}
\[
(s\mI - \mA)^{-1} = \begin{bmatrix} s+1 & 0 \\ 0 & s+2 \end{bmatrix}^{-1} = \begin{bmatrix} \frac{1}{s+1} & 0 \\ 0 & \frac{1}{s+2} \end{bmatrix}
\]

Taking the inverse Laplace transform:
\[
\ee^{\mA t} = \iLaplace\left\{\begin{bmatrix} \frac{1}{s+1} & 0 \\ 0 & \frac{1}{s+2} \end{bmatrix}\right\} = \begin{bmatrix} \ee^{-t} & 0 \\ 0 & \ee^{-2t} \end{bmatrix}
\]

\textbf{Interpretation:} The state transition matrix $\ee^{\mA t}$ shows that both modes decay exponentially, with eigenvalues $\lambda_1 = -1$ and $\lambda_2 = -2$ determining the decay rates. The system is stable since both eigenvalues have negative real parts.
\end{examplebox}

\section{Linear Systems of Equations}

Solving $\mA\vecx = \vecb$:
\begin{itemize}
    \item \textbf{Unique solution} if $\mA$ is square and invertible: $\vecx = \mA^{-1}\vecb$
    \item \textbf{No solution} if $\vecb \notin \range(\mA)$
    \item \textbf{Infinite solutions} if $\nullspace(\mA) \neq \{\mathbf{0}\}$
\end{itemize}

\begin{examplebox}[title=Example: Solving a System Using Gaussian Elimination]
Solve the system $\mA\vecx = \vecb$ where:
\[
\begin{bmatrix} 2 & 1 & -1 \\ -3 & -1 & 2 \\ -2 & 1 & 2 \end{bmatrix}\begin{bmatrix} x_1 \\ x_2 \\ x_3 \end{bmatrix} = \begin{bmatrix} 8 \\ -11 \\ -3 \end{bmatrix}
\]

\textbf{Solution:}

\textbf{Step 1: Form augmented matrix}
\[
\left[\begin{array}{ccc|c} 2 & 1 & -1 & 8 \\ -3 & -1 & 2 & -11 \\ -2 & 1 & 2 & -3 \end{array}\right]
\]

\textbf{Step 2: Forward elimination}

$R_2 \leftarrow R_2 + \frac{3}{2}R_1$:
\[
\left[\begin{array}{ccc|c} 2 & 1 & -1 & 8 \\ 0 & 1/2 & 1/2 & 1 \\ -2 & 1 & 2 & -3 \end{array}\right]
\]

$R_3 \leftarrow R_3 + R_1$:
\[
\left[\begin{array}{ccc|c} 2 & 1 & -1 & 8 \\ 0 & 1/2 & 1/2 & 1 \\ 0 & 2 & 1 & 5 \end{array}\right]
\]

$R_3 \leftarrow R_3 - 4R_2$:
\[
\left[\begin{array}{ccc|c} 2 & 1 & -1 & 8 \\ 0 & 1/2 & 1/2 & 1 \\ 0 & 0 & -1 & 1 \end{array}\right]
\]

\textbf{Step 3: Back substitution}

From Row 3: $-x_3 = 1 \Rightarrow x_3 = -1$

From Row 2: $\frac{1}{2}x_2 + \frac{1}{2}(-1) = 1 \Rightarrow x_2 = 3$

From Row 1: $2x_1 + 3 - (-1) = 8 \Rightarrow 2x_1 = 4 \Rightarrow x_1 = 2$

\textbf{Solution:} $\vecx = \begin{bmatrix} 2 \\ 3 \\ -1 \end{bmatrix}$

\textbf{Verification:} $\mA\vecx = \begin{bmatrix} 4+3+1 \\ -6-3-2 \\ -4+3-2 \end{bmatrix} = \begin{bmatrix} 8 \\ -11 \\ -3 \end{bmatrix} = \vecb$ \checkmark
\end{examplebox}

\subsubsection{Least Squares}
When $\mA\vecx = \vecb$ has no exact solution:
\begin{equation}
\hat{\vecx} = (\mA^\top\mA)^{-1}\mA^\top\vecb = \mA^\dagger\vecb
\end{equation}
where $\mA^\dagger$ is the Moore-Penrose pseudoinverse.

\section{Controllability and Observability}

Two fundamental concepts in state-space control theory determine whether a system can be fully controlled and whether its internal states can be determined from measurements.

\begin{definitionbox}[title=Definition 0.7: Controllability Matrix]
For a linear time-invariant system $\dot{\vecx} = \mA\vecx + \mB\vecu$, the \textbf{controllability matrix} is:
\begin{equation}
\mathcal{C} = \begin{bmatrix} \mB & \mA\mB & \mA^2\mB & \cdots & \mA^{n-1}\mB \end{bmatrix}
\end{equation}
where $n$ is the dimension of the state vector $\vecx$.
\end{definitionbox}

\begin{keyconceptbox}[title=Controllability Condition]
A system is \textbf{controllable} (or reachable) if and only if the controllability matrix has full row rank:
\[
\rank(\mathcal{C}) = n
\]

Physical interpretation: The system is controllable if the input $\vecu$ can drive the state from any initial condition to any desired final state in finite time.
\end{keyconceptbox}

\begin{definitionbox}[title=Definition 0.8: Observability Matrix]
For a linear time-invariant system with output $\vecy = \mC\vecx$, the \textbf{observability matrix} is:
\begin{equation}
\mathcal{O} = \begin{bmatrix} \mC \\ \mC\mA \\ \mC\mA^2 \\ \vdots \\ \mC\mA^{n-1} \end{bmatrix}
\end{equation}
\end{definitionbox}

\begin{keyconceptbox}[title=Observability Condition]
A system is \textbf{observable} if and only if the observability matrix has full column rank:
\[
\rank(\mathcal{O}) = n
\]

Physical interpretation: The system is observable if the initial state $\vecx(0)$ can be uniquely determined from knowledge of the input $\vecu(t)$ and output $\vecy(t)$ over a finite time interval.
\end{keyconceptbox}

\subsubsection{Connection to State-Space Control}

\begin{warningbox}[title=Importance for Control Design]
\begin{itemize}
    \item \textbf{Pole placement} (state feedback): Requires controllability. If the system is not controllable, some eigenvalues cannot be moved by state feedback.
    \item \textbf{Observer design}: Requires observability. If the system is not observable, some states cannot be estimated from output measurements.
    \item \textbf{Kalman decomposition}: Any LTI system can be decomposed into controllable/uncontrollable and observable/unobservable subsystems.
\end{itemize}
\end{warningbox}

The duality between controllability and observability is a fundamental result:
\begin{itemize}
    \item The pair $(\mA, \mB)$ is controllable if and only if $(\mA^\top, \mB^\top)$ is observable.
    \item The pair $(\mA, \mC)$ is observable if and only if $(\mA^\top, \mC^\top)$ is controllable.
\end{itemize}

This duality is essential for understanding the relationship between controller design and observer design in modern control theory.

% ============================================================================
